\section{求解代码清单}

本文全部求解代码位于项目 \texttt{code/} 目录，每个问题一个独立脚本，按序运行即可复现全文结果。所有代码均为 Python 3 实现，依赖标准科学计算与绘图库（numpy、pandas、matplotlib、seaborn、networkx，以及用于整数规划与二分求解的 scipy），不含随机种子依赖，重复运行结果完全一致。

% 按《人工智能工具及输出使用规定》第 5 条，使用人工智能工具辅助编程的，应在程序前添加
% 注释说明工具名称、版本/型号、开发机构与版本发布日期；各 .py 文件首部均已按要求标注。
% 附录内的表号需与「附录 A」的字母编号一致，这里改用不浮动表格，避免跨节漂移。
\begin{center}
\begin{minipage}{\textwidth}
	\captionof{table}{求解代码与结果文件清单}
	\label{tab:code-list}
	\centering
	\begin{tabularx}{0.98\textwidth}{lX}
		\toprule
		文件 & 说明 \\
		\midrule
		code/core.py & 公共物理模型：数据加载、DEM 采样与遮挡判定、航段时间/能耗/航程、充电模型、通信链路模型 \\
		code/q1.py & 问题一：最大安全载荷二分求解、同质类多重集动态规划组批、指派式与集合划分整数规划双路径复核、$\rho$--$\kappa$ 二维灵敏度 \\
		code/q2.py & 问题二：硬时限约束下的架次构造与共享电池离散事件调度、ALNS 破坏--修复搜索、$\varepsilon$-约束扫描、消融实验与记忆化 \\
		code/q2\_exact.py & 问题二的两层小规模精确验证：列生成＋集合划分整数规划（分区层）、big-$M$ 析取混合整数规划（时序层） \\
		code/q3.py & 问题三：轨迹采样与直连失效区间提取、三维悬停站集合覆盖整数规划、运输错峰与中继架次调度的联合优化 \\
		code/q4.py & 问题四：原子任务单元并查集归并、受限增长串完全枚举、组列集合划分整数规划复核、分组峰值资源需求与缺口核算 \\
		code/verify.py & 独立复核脚本：只读结果表重算第 9 章的全部硬约束、通信连续性与 DEM 离散化误差 \\
		code/fill\_results.py & 把 \texttt{results/*.csv} 回填进《结果提交模板》生成上传用结果表 \\
		code/figures.py & 数据类配图生成脚本（地形图、灵敏度图、直连失效区分布图） \\
		code/figures\_adv.py & 多面板高级配图脚本（帕累托前沿、ALNS 诊断、通信状态时序带、资源雷达图等） \\
		code/make\_roadmap.py & 生成总技术路线图的 .drawio 源文件（版式与图内数值的来源） \\
		code/render\_drawio.py & 把 figures/ 下的 .drawio 示意图离线渲染为 PNG/PDF \\
		results/*.csv & 各问题结果表（对应结果提交模板各 sheet） \\
		figures/*.drawio & 示意图的可编辑源文件（总技术路线图、问题二算法流程） \\
		figures/*.png & 论文插图 \\
		\bottomrule
	\end{tabularx}
\end{minipage}
\end{center}

\section{主要结果数据}

% 依《人工智能工具及输出使用规定》第 4 条，在数据分析结果处加注工具信息。
% 两段都含长串不可断的文件名与英文模型名，必须一起放进 sloppypar，
% 否则 TeX 无法找到合法断点，会产生 Overfull \hbox。
\begin{sloppypar}
\noindent\textbf{数据来源注释：}本文全部数据分析结果由本队在人工智能工具辅助下编写的求解脚本计算得到，所用工具为 DeepSeek，版本/型号 DeepSeek V4.1 Flash（API 模型名 \texttt{deepseek-flash}），版本发布日期 2026 年 9 月 10 日，开发机构为杭州深度求索人工智能基础技术研究有限公司；各脚本首部与支撑材料《AI 工具使用详情》中已列明工具信息与使用过程。全部数据均由本队独立重复运行核验。

问题一最大安全载荷、推荐组批方案、帕累托前沿与双路径复核结果见 \texttt{results/q1\_max\_payload.csv}、\texttt{q1\_recommended\_batching.csv}、\texttt{q1\_pareto.csv}、\texttt{q1\_ilp\_gap.csv} 与 \texttt{q1\_sensitivity.csv}；问题二运输架次、逐箱交付时刻、推荐方案快照、非支配前沿、消融与精确验证结果见 \texttt{results/q2\_transport\_trips.csv}、\texttt{q2\_box\_delivery.csv}、\texttt{q2\_recommended.json}、\texttt{q2\_pareto.csv}、\texttt{q2\_ablation.csv} 与 \texttt{q2\_exact.csv}；问题三悬停站选址、失效区间绑定、中继架次、错峰量、通信阶段与三种口径的覆盖率对照见 \texttt{results/q3\_stations.csv}、\texttt{q3\_intervals.csv}、\texttt{q3\_relay\_trips.csv}、\texttt{q3\_stagger.csv}、\texttt{q3\_comm\_phases.csv} 与 \texttt{q3\_coverage.csv}；问题四原子任务单元、分区配置、全部分区枚举与两类方案对照见 \texttt{results/q4\_units.csv}、\texttt{q4\_partition.csv}、\texttt{q4\_all\_partitions.csv} 与 \texttt{q4\_comparison.csv}。
\end{sloppypar}
